% GENERATED FILE. Edit canonical lessons, then run npm run generate:books.
% canonical-source-hash: fnv1a64:5449bcaf6168b44e
% canonical-lessons: GU-C167-bhido, GU-C167-vatano, GU-C167-idu, GU-C167-baphvu, GU-C167-talvu-fry

\chapter{Okra, peas and eggs}
\label{ch:gu-okra-peas-and-eggs}
\hlchaptermodality{167}

\begin{chapteropening}
\textbf{By the end of this chapter:} \emph{I can talk about vegetables and eggs and how they are cooked.}

Complete the last lesson of chapter 167: I can talk about vegetables and eggs and how they are cooked.
\end{chapteropening}

\section[\texorpdfstring{\gu{ભીંડો} (bhī̃ḍo)}{ભીંડો (bhī̃ḍo)}]{\gu{ભીંડો} (bhī̃ḍo) — okra, lady's finger}

\label{lesson:GU-C167-bhido}



\begin{quote}
Before the new one: say the Gujarati for fresh coriander, then the Gujarati for to grind.
\end{quote}

\subsection*{You'll want to know: \gu{ભીંડો}}
\textbf{\gu{ભીંડો}} — \emph{bhī̃ḍo} — \textquotedblleft{}okra, lady's finger\textquotedblright{}.

\begin{quote}
\textbf{\gu{મને} \gu{ભીંડાનું} \gu{શાક} \gu{ગમે} \gu{છે}.} — \emph{mane bhī̃ḍānũ śāk game che} — \textquotedblleft{}I like okra curry.\textquotedblright{}
\end{quote}

\begin{cousinweb}
\textbf{\gu{ભીંડો}} (\emph{bhī̃ḍo}) is a tadbhava. It comes from Sanskrit \emph{bhiṇḍā}, the vegetable \emph{Abelmoschus esculentus}, and CDIAL lists the Gujarati \emph{bhĩḍo} (CDIAL 9492).
\end{cousinweb}

\begin{grammarlens}[title={What you've built}]
One more word for everyday life at A2.
\end{grammarlens}

\subsection*{Your turn}
\begin{itemize}
\raggedright
  \item \emph{Say it:} \emph{bhī̃ḍo}
  \item \emph{Say it:} \emph{bhī̃ḍo}, once more
  \item \emph{Say it:} \emph{mane bhī̃ḍānũ śāk game che} — I like okra curry.
  \item \emph{Recall:} say the Gujarati for fresh coriander, then the Gujarati for to grind, then say \emph{bhī̃ḍo} again
\end{itemize}

\begin{tcolorbox}[breakable,colback=teal!4,colframe=teal!35!black,title={Before you move on}]
What does \gu{ભીંડો} mean? (\textquotedblleft{}Okra, lady's finger\textquotedblright{}.) Say it once more.
\end{tcolorbox}
\section[\texorpdfstring{\gu{વટાણો} (vaṭāṇo)}{વટાણો (vaṭāṇo)}]{\gu{વટાણો} (vaṭāṇo) — pea}

\label{lesson:GU-C167-vatano}



\begin{quote}
Before the new one: say the Gujarati for to grind, then the Gujarati for okra.
\end{quote}

\subsection*{You'll want to know: \gu{વટાણો}}
\textbf{\gu{વટાણો}} — \emph{vaṭāṇo} — \textquotedblleft{}pea\textquotedblright{}.

\begin{quote}
\textbf{\gu{બાળક} \gu{વટાણા} \gu{ગણે} \gu{છે}.} — \emph{bāḷak vaṭāṇā gaṇe che} — \textquotedblleft{}The child counts the peas.\textquotedblright{}
\end{quote}

Mostly in the plural, \textbf{\gu{વટાણા}}, \emph{vaṭāṇā}.

\begin{cousinweb}
The origin of \textbf{\gu{વટાણો}} (\emph{vaṭāṇo}) is not settled. CDIAL lists no Gujarati form of it, and EtymDB has no entry for it.
\end{cousinweb}

\begin{grammarlens}[title={What you've built}]
One more word for everyday life at A2.
\end{grammarlens}

\subsection*{Your turn}
\begin{itemize}
\raggedright
  \item \emph{Say it:} \emph{vaṭāṇo}
  \item \emph{Say it:} \emph{vaṭāṇo}, once more
  \item \emph{Say it:} \emph{bāḷak vaṭāṇā gaṇe che} — The child counts the peas.
  \item \emph{Recall:} say the Gujarati for to grind, then the Gujarati for okra, then say \emph{vaṭāṇo} again
\end{itemize}

\begin{tcolorbox}[breakable,colback=teal!4,colframe=teal!35!black,title={Before you move on}]
What does \gu{વટાણો} mean? (\textquotedblleft{}Pea\textquotedblright{}.) Say it once more.
\end{tcolorbox}
\section[\texorpdfstring{\gu{ઈંડું} (ī̃ḍũ)}{ઈંડું (ī̃ḍũ)}]{\gu{ઈંડું} (ī̃ḍũ) — egg}

\label{lesson:GU-C167-idu}



\begin{quote}
Before the new one: say the Gujarati for okra, then the Gujarati for pea.
\end{quote}

\subsection*{You'll want to know: \gu{ઈંડું}}
\textbf{\gu{ઈંડું}} — \emph{ī̃ḍũ} — \textquotedblleft{}egg\textquotedblright{}.

\begin{quote}
\textbf{\gu{તે} \gu{નાસ્તામાં} \gu{ઈંડું} \gu{ખાય} \gu{છે}.} — \emph{te nāstāmā̃ ī̃ḍũ khāy che} — \textquotedblleft{}He eats an egg for breakfast.\textquotedblright{}
\end{quote}

\begin{cousinweb}
\textbf{\gu{ઈંડું}} (\emph{ī̃ḍũ}) is a tadbhava. It comes from Sanskrit \emph{āṇḍa}, \textquotedblleft{}egg\textquotedblright{}, and CDIAL lists the Gujarati \emph{ĩḍũ}, \textquotedblleft{}egg\textquotedblright{} (CDIAL 1111).
\end{cousinweb}

\begin{grammarlens}[title={What you've built}]
One more word for everyday life at A2.
\end{grammarlens}

\subsection*{Your turn}
\begin{itemize}
\raggedright
  \item \emph{Say it:} \emph{ī̃ḍũ}
  \item \emph{Say it:} \emph{ī̃ḍũ}, once more
  \item \emph{Say it:} \emph{te nāstāmā̃ ī̃ḍũ khāy che} — He eats an egg for breakfast.
  \item \emph{Recall:} say the Gujarati for okra, then the Gujarati for pea, then say \emph{ī̃ḍũ} again
\end{itemize}

\begin{tcolorbox}[breakable,colback=teal!4,colframe=teal!35!black,title={Before you move on}]
What does \gu{ઈંડું} mean? (\textquotedblleft{}Egg\textquotedblright{}.) Say it once more.
\end{tcolorbox}
\section[\texorpdfstring{\gu{બાફવું} (bāphvũ)}{બાફવું (bāphvũ)}]{\gu{બાફવું} (bāphvũ) — to steam, to boil soft}

\label{lesson:GU-C167-baphvu}



\begin{quote}
Before the new one: say the Gujarati for pea, then the Gujarati for egg.
\end{quote}

\subsection*{You'll want to know: \gu{બાફવું}}
\textbf{\gu{બાફવું}} — \emph{bāphvũ} — \textquotedblleft{}to steam, to boil soft\textquotedblright{}.

\begin{quote}
\textbf{\gu{તે} \gu{શાક} \gu{બાફે} \gu{છે}.} — \emph{te śāk bāphe che} — \textquotedblleft{}She steams the vegetables.\textquotedblright{}
\end{quote}

\begin{cousinweb}
\textbf{\gu{બાફવું}} (\emph{bāphvũ}) is a tadbhava. It comes from Sanskrit \emph{bāṣpāyate}, \textquotedblleft{}emits vapour\textquotedblright{}, built on \emph{bāṣpa}, \textquotedblleft{}steam, tears\textquotedblright{}, and CDIAL glosses the Gujarati verb \textquotedblleft{}to boil\textquotedblright{} (CDIAL 9225).
\end{cousinweb}

\begin{grammarlens}[title={What you've built}]
One more word for everyday life at A2.
\end{grammarlens}

\subsection*{Your turn}
\begin{itemize}
\raggedright
  \item \emph{Say it:} \emph{bāphvũ}
  \item \emph{Say it:} \emph{bāphvũ}, once more
  \item \emph{Say it:} \emph{te śāk bāphe che} — She steams the vegetables.
  \item \emph{Recall:} say the Gujarati for pea, then the Gujarati for egg, then say \emph{bāphvũ} again
\end{itemize}

\begin{tcolorbox}[breakable,colback=teal!4,colframe=teal!35!black,title={Before you move on}]
What does \gu{બાફવું} mean? (\textquotedblleft{}To steam, to boil soft\textquotedblright{}.) Say it once more.
\end{tcolorbox}
\section[\texorpdfstring{\gu{તળવું} (taḷvũ)}{તળવું (taḷvũ)}]{\gu{તળવું} (taḷvũ) — to fry}

\label{lesson:GU-C167-talvu-fry}



\begin{quote}
Before the new one: say the Gujarati for egg, then the Gujarati for to steam.
\end{quote}

\subsection*{You'll want to know: \gu{તળવું}}
\textbf{\gu{તળવું}} — \emph{taḷvũ} — \textquotedblleft{}to fry\textquotedblright{}.

\begin{quote}
\textbf{\gu{માતા} \gu{તેલમાં} \gu{ભજિયાં} \gu{તળે} \gu{છે}.} — \emph{mātā telmā̃ bhajiyā̃ taḷe che} — \textquotedblleft{}Mother fries bhajiyas in oil.\textquotedblright{}
\end{quote}

\begin{cousinweb}
\textbf{\gu{તળવું}} (\emph{taḷvũ}) is a tadbhava. It comes from \emph{*talati}, \textquotedblleft{}fries\textquotedblright{}, a reconstructed form beside Sanskrit \emph{talita}, \textquotedblleft{}fried\textquotedblright{}, and CDIAL glosses the Gujarati verb \textquotedblleft{}to fry, cook in ghee\textquotedblright{} (CDIAL 5736).
\end{cousinweb}

\begin{grammarlens}[title={What you've built}]
One more word for everyday life at A2.
\end{grammarlens}

\subsection*{Your turn}
\begin{itemize}
\raggedright
  \item \emph{Say it:} \emph{taḷvũ}
  \item \emph{Say it:} \emph{taḷvũ}, once more
  \item \emph{Say it:} \emph{mātā telmā̃ bhajiyā̃ taḷe che} — Mother fries bhajiyas in oil.
  \item \emph{Recall:} say the Gujarati for egg, then the Gujarati for to steam, then say \emph{taḷvũ} again
\end{itemize}

\begin{tcolorbox}[breakable,colback=teal!4,colframe=teal!35!black,title={Before you move on}]
What does \gu{તળવું} mean? (\textquotedblleft{}To fry\textquotedblright{}.) Say it once more.
\end{tcolorbox}
